\section{The parameterized semiprime denominator system}
\label{sec:denominator-system}

Fix throughout a reduced target
\[
 t=\frac ab\in(0,1)
\]
with squarefree $b$, a fixed anchor ratio $\eta\in(0,1)$, a finite forbidden set $\mathcal T$, and a fixed exponent $\gamma>1$.  Reducedness and $0<t<1$ imply $b>1$.  Thus the set $S_b$ of prime divisors of $b$ is nonempty and
\begin{equation}
\label{eq:tau-positive}
 \tau(b):=\sum_{r\in S_b}\frac1{r^2}>0.
\end{equation}
Put
\begin{equation}
\label{eq:lambda-gamma}
 \lambda_\gamma=\frac{(\log\gamma)^2}{2}.
\end{equation}
We assume the exact inequalities for the centred no-wrap construction
\begin{equation}
\label{eq:admissible-region}
 t<\lambda_\gamma<1.
\end{equation}
Equivalently,
\begin{equation}
\label{eq:admissible-gamma}
 \exp(\sqrt{2t})<\gamma<\exp(\sqrt2).
\end{equation}
All asymptotic statements are for $X\to\infty$ with $t$, $\gamma$, $\eta$, $b$ and $\mathcal T$ fixed.  Unless stated otherwise, implied constants may depend on the fixed tuple $(t,\gamma,\eta,b)$ and on fixed numerical choices, but not on $X$ or on the later Gaussian cutoff $C$.  The estimates below are uniform when $\eta$ ranges in a fixed compact subset of $(0,1)$.

\subsection{Prime blocks and their power sums}

Put
\begin{equation}
\label{eq:prime-blocks}
 Z=X^\gamma,
 \qquad
 P=\{p\text{ prime}:X\leq p<Z\},
 \qquad
 B=\{q\text{ prime}:\eta Z\leq q<Z\},
\end{equation}
and let $S_b$ be the set of prime divisors of $b$.  We take $X$ so large that every prime in $P$ exceeds every prime divisor of $b$, and every denominator constructed below exceeds every element of $\mathcal T$.

We use the prime number theorem in the forms
\[
 \pi(y)\sim\frac{y}{\log y},
 \qquad
 \vartheta(y):=\sum_{p\leq y}\log p=y+o(y),
\]
together with their standard partial-summation consequences; see Montgomery and Vaughan~\cite[Chapter~1]{montgomery-vaughan-2007}.

\begin{lemma}[Prime supply in fixed-ratio intervals]
\label{lem:fixed-ratio-prime-supply}
For fixed $0<u<v$, as $y\to\infty$,
\[
 \#\{p\text{ prime}:uy\leq p<vy\}
 =\bigl(v-u+o_{u,v}(1)\bigr)\frac{y}{\log y}.
\]
The error is uniform when $(u,v)$ ranges over a compact subset of $\{(u,v):0<u<v\}$.
\end{lemma}

\begin{proof}
Apply the prime number theorem at $uy$ and $vy$.  Replacing $\log(uy)$ and $\log(vy)$ by $\log y$ introduces only a relative $o_{u,v}(1)$, uniformly on compact parameter sets.
\end{proof}

Define
\[
 S_X=\sum_{p\in P}\frac1p,
 \qquad
 U_X=\sum_{p\in P}\frac1{p^2},
 \qquad
 M_B=|B|,
\]
and
\[
 V_B=\sum_{q\in B}\frac1q,
 \qquad
 W_B=\sum_{q\in B}\frac1{q^2}.
\]

\begin{lemma}[Parameterized prime reciprocal estimates]
\label{lem:reciprocal-prime-estimates}
As $X\to\infty$,
\begin{align}
 S_X&=\log\gamma+o(1),
 \label{eq:SX-asymptotic}\\
 U_X&\sim\frac1{X\log X},
 \label{eq:UX-asymptotic}\\
 |P|&\sim\frac{Z}{\log Z},
 \label{eq:P-cardinality}\\
 M_B&=\bigl(1-\eta+o_\eta(1)\bigr)\frac{Z}{\log Z},
 \label{eq:B-cardinality}\\
 V_B&=\frac{\log(1/\eta)+o_\eta(1)}{\log Z},
 \label{eq:B-reciprocal-mass}\\
 W_B&=\frac{\eta^{-1}-1+o_\eta(1)}{Z\log Z}.
 \label{eq:B-power-sums}
\end{align}
\end{lemma}

\begin{proof}
Abel summation gives
\[
 \sum_{X\leq p<X^\gamma}\frac1p
 =\int_X^{X^\gamma}\frac{du}{u\log u}+o(1)
 =\log\gamma+o(1).
\]
A second partial summation gives
\[
 \sum_{X\leq p<X^\gamma}\frac1{p^2}
 =\int_X^{X^\gamma}\frac{du}{u^2\log u}(1+o(1))
 \sim\frac1{X\log X};
\]
the upper endpoint is negligible because $\gamma>1$ is fixed.  The cardinality statements follow from \cref{lem:fixed-ratio-prime-supply}.  Partial summation on $[\eta Z,Z)$ gives
\[
 \sum_{\eta Z\leq q<Z}\frac1q
 =\int_{\eta Z}^{Z}\frac{du}{u\log u}+o_\eta(1/\log Z)
 =\frac{\log(1/\eta)+o_\eta(1)}{\log Z}
\]
and
\[
 \sum_{\eta Z\leq q<Z}\frac1{q^2}
 =\int_{\eta Z}^{Z}\frac{du}{u^2\log u}(1+o_\eta(1))
 =\frac{\eta^{-1}-1+o_\eta(1)}{Z\log Z}.
\]
All fixed-ratio estimates are uniform for $\eta$ in compact subsets of $(0,1)$.
\end{proof}

\subsection{The complete pair family and target rows}

Define
\[
 \mathcal E_{\mathrm{pair}}
 =\{pq:p,q\in P,\ p<q\},
 \qquad
 \mathcal E_b
 =\{rq:r\in S_b,\ q\in B\},
\]
and
\begin{equation}
\label{eq:E-and-L}
 \mathcal E=\mathcal E_{\mathrm{pair}}\sqcup\mathcal E_b,
 \qquad
 L=b\prod_{p\in P}p.
\end{equation}
The disjoint union follows from unique factorization.  Put
\[
 \Lambda=\sum_{e\in\mathcal E}\frac1e.
\]

\begin{proposition}[Parameterized arithmetic capacity]
\label{prop:arithmetic-capacity}
For all sufficiently large $X$:
\begin{enumerate}[label=(\roman*)]
\item every $e\in\mathcal E$ is a squarefree semiprime;
\item the elements of $\mathcal E$ are pairwise distinct, avoid $\mathcal T$, and divide $L$;
\item one has
\begin{equation}
\label{eq:load-limit}
 \Lambda\longrightarrow\lambda_\gamma;
\end{equation}
more precisely, for fixed constants
\[
 t<\lambda_-<\lambda_\gamma<\lambda_+<1,
\]
one has $\lambda_-<\Lambda<\lambda_+$ for all sufficiently large $X$;
\item the complete square load satisfies
\begin{equation}
\label{eq:complete-square-load}
 \sum_{e\in\mathcal E}\frac1{e^2}
 =\frac{1+o(1)}{2X^2\log^2X}
  +\frac{(\eta^{-1}-1)\tau(b)+o_{b,\eta}(1)}{Z\log Z}.
\end{equation}
\end{enumerate}
\end{proposition}

\begin{proof}
Only the load and square load require proof.  The exact pair identity \cref{cor:pair-identity} gives
\begin{equation}
\label{eq:pair-load}
 \Lambda_{\mathrm{pair}}
 =\frac12(S_X^2-U_X)
 \longrightarrow\frac{(\log\gamma)^2}{2}=\lambda_\gamma.
\end{equation}
The target-row load is
\[
 \Lambda_b
 =\left(\sum_{r\in S_b}\frac1r\right)V_B
 =\left(\sum_{r\in S_b}\frac1r\right)
   \frac{\log(1/\eta)+o_\eta(1)}{\log Z}
 =o_{b,\eta}(1),
\]
which proves \cref{eq:load-limit} and the fixed load interval.

For the pair square load, apply the pair identity to $u_p=p^{-2}$:
\[
 \sum_{p<q}\frac1{p^2q^2}
 =\frac12\left(U_X^2-\sum_{p\in P}\frac1{p^4}\right)
 \sim\frac1{2X^2\log^2X}.
\]
Indeed the fourth-power diagonal is $O(1/(X^3\log X))$.  The target rows contribute
\[
 \sum_{r\in S_b}\sum_{q\in B}\frac1{r^2q^2}
 =\tau(b)W_B
 =\frac{(\eta^{-1}-1)\tau(b)+o_{b,\eta}(1)}{Z\log Z}.
\]
This proves \cref{eq:complete-square-load} without assuming which summand dominates.
\end{proof}

\begin{proposition}[Family size and modulus]
\label{prop:family-size-modulus}
One has
\begin{equation}
\label{eq:family-size-modulus}
 M_X:=|\mathcal E|
 =(1+o(1))\frac{Z^2}{2\log^2Z},
 \qquad
 \log L=(1+o(1))Z=o(M_X).
\end{equation}
The target rows have
\[
 |\mathcal E_b|
 =\bigl(|S_b|(1-\eta)+o_{b,\eta}(1)\bigr)\frac{Z}{\log Z}
 =O_{b,\eta}(Z/\log Z)
\]
elements and do not affect the leading family size.  Every denominator is at most $Z^2$ for large $X$.
\end{proposition}

\begin{proof}
The pair family has $\binom{|P|}{2}=(1+o(1))Z^2/(2\log^2Z)$ elements by \cref{eq:P-cardinality}, while the displayed target-row count follows from \cref{eq:B-cardinality}.  This proves the formula for $M_X$.

For the modulus, the prime number theorem in the form $\vartheta(y)=y+o(y)$ gives
\begin{align*}
 \log L
 &=\log b+\sum_{p\in P}\log p\\
 &=\vartheta(Z)-\vartheta(X)+O_b(1)\\
 &=(1+o(1))Z.
\end{align*}
The harmless $O_b(1)$ accounts for the endpoint convention and the fixed factor $b$.  Since $\gamma>1$, one has $X=o(Z)$.  Finally,
\[
 \frac{\log L}{M_X}
 =(2+o(1))\frac{\log^2Z}{Z}\longrightarrow0,
\]
so $\log L=o(M_X)$.
\end{proof}

\subsection{Exact centring and the variance boundary}

Set
\begin{equation}
\label{eq:theta-definition}
 \theta=\frac{t}{\Lambda}.
\end{equation}
By \cref{prop:arithmetic-capacity}, $\theta$ lies in a fixed compact interval $I_{t,\gamma,\eta}\subset(0,1)$ and
\begin{equation}
\label{eq:theta-limit}
 \theta\longrightarrow
 \alpha_{t,\gamma}
 :=\frac{t}{\lambda_\gamma}
 =\frac{2t}{(\log\gamma)^2}.
\end{equation}
Moreover,
\begin{equation}
\label{eq:exact-centering-load}
 \theta\sum_{e\in\mathcal E}\frac1e=t,
 \qquad
 \Lambda<1.
\end{equation}
Define the actual-family variance by
\begin{equation}
\label{eq:sigma-E}
 \sigma_{\mathcal E}^2
 =\theta(1-\theta)\sum_{e\in\mathcal E}\frac1{e^2}.
\end{equation}

\begin{proposition}[Total actual-family variance]
\label{prop:variance-regimes}
Let $\alpha=\alpha_{t,\gamma}$.  Then
\begin{equation}
\label{eq:total-variance}
 \sigma_{\mathcal E}^2
 \sim \alpha(1-\alpha)
 \left(
  \frac1{2X^2\log^2X}
  +\frac{(\eta^{-1}-1)\tau(b)}{Z\log Z}
 \right).
\end{equation}
\end{proposition}

\begin{proof}
Combine \cref{eq:theta-limit,eq:complete-square-load,eq:sigma-E}.  Both summands are retained, so the statement is valid without first deciding which contribution dominates.
\end{proof}

\begin{corollary}[Explicit variance regimes]
\label{cor:variance-regimes}
With $\alpha=\alpha_{t,\gamma}$:
\begin{enumerate}[label=(\Alph*)]
\item If $\gamma>2$, the pair square load dominates and
\begin{equation}
\label{eq:variance-gamma-large}
 \sigma_{\mathcal E}^2
 \sim\frac{\alpha(1-\alpha)}{2X^2\log^2X}.
\end{equation}
\item If $\gamma=2$, the target-row square load dominates the pair square load by a factor of order $\log X$, and
\begin{equation}
\label{eq:variance-gamma-two}
 \sigma_{\mathcal E}^2
 \sim\frac{\alpha(1-\alpha)(\eta^{-1}-1)\tau(b)}{Z\log Z}.
\end{equation}
\item If $1<\gamma<2$, the target-row square load dominates by a power of $X$, and the same asymptotic \cref{eq:variance-gamma-two} holds.
\end{enumerate}
\end{corollary}

\begin{proof}
The ratio of the target-row summand in \cref{eq:total-variance} to the pair summand is asymptotic, up to a fixed positive factor depending on $(\eta,b)$, to
\[
 \frac{X^2\log^2X}{Z\log Z}
 \asymp X^{2-\gamma}\log X.
\]
It tends to $0$ for $\gamma>2$, to infinity logarithmically for $\gamma=2$, and to infinity polynomially for $\gamma<2$.  Positivity of $(\eta^{-1}-1)\tau(b)$ follows from $\eta\in(0,1)$ and \cref{eq:tau-positive}; the case $b=1$ cannot occur in the direct setting $t\in(0,1)$.
\end{proof}

Thus $\gamma=2$ is a genuine variance boundary for every fixed $\eta\in(0,1)$, but it does not require a separate proof architecture: every later estimate is stated first with the total variance \cref{eq:total-variance}.

The finite Fourier numerator is
\begin{equation}
\label{eq:Fh-definition}
 F(h)=\exp\left(-\frac{2\pi i a h}{b}\right)
 \prod_{e\in\mathcal E}
 \left((1-\theta)+\theta\exp\left(\frac{2\pi i h}{e}\right)\right).
\end{equation}
Compactness of $I_{t,\gamma,\eta}$ gives a constant
\[
 \kappa_{\rm mod}=\kappa_{\rm mod}(t,\gamma,\eta)>0
\]
such that
\begin{equation}
\label{eq:bernoulli-modulus}
 |(1-\theta)+\theta\exp(2\pi i u)|
 \leq\exp(-\kappa_{\rm mod}\dist{u}^2).
\end{equation}
Consequently,
\begin{equation}
\label{eq:global-Fourier-energy}
 |F(h)|\leq\exp(-\kappa_{\rm mod}Q_{\mathcal E}(h)),
 \qquad
 Q_{\mathcal E}(h)=\sum_{e\in\mathcal E}\dist{h/e}^2.
\end{equation}

\subsection{Exact factor partition}

Because $b$ and $\prod_{p\in P}p$ are coprime and squarefree, the period $L$ is squarefree.  Chinese-remainder coordinates split the frequency into those indexed by $B$, by $P\setminus B$, and by $S_b$.  The denominator factors are assigned exactly once as follows:
\begin{center}
\begin{tabular}{ll}
\toprule
Denominator class & role \\
\midrule
$q q'$ with $q,q'\in B$ & internal anchor factor \\
$r q$ with $r\in P\setminus B$, $q\in B$ & lower-prime anchor-cross row $r$ \\
$r q$ with $r\in S_b$, $q\in B$ & target-coordinate row $r$ \\
$p p'$ with $p,p'\in P\setminus B$ & retained skeleton factor \\
\bottomrule
\end{tabular}
\end{center}
The target character and every phase not absorbed into a row remain in the complex residual factor.  The four denominator classes are disjoint and cover $\mathcal E$.
